% ============================================================================
%  Multiphysics Hub -- technical report
%  B.Sc. Computational Engineering Science (CES), RWTH Aachen
%
%  COMPILE: upload this folder (report.tex + figures/) to Overleaf and press
%  "Recompile" (pdfLaTeX). Locally: pdflatex report.tex (run twice).
% ============================================================================
\documentclass[11pt,a4paper]{article}

\usepackage[T1]{fontenc}
\usepackage[utf8]{inputenc}
\usepackage{lmodern}
\usepackage[margin=2.3cm]{geometry}
\usepackage{amsmath,amssymb}
\usepackage{graphicx}
\usepackage{booktabs}
\usepackage{siunitx}
\usepackage{caption}
\usepackage{subcaption}
\usepackage[hidelinks]{hyperref}
\usepackage{xcolor}
\usepackage{enumitem}

\sisetup{detect-all}
\graphicspath{{figures/}}
\captionsetup{font=small,labelfont=bf}
\setlist{nosep,leftmargin=1.4em}

% --- RWTH house colours -----------------------------------------------------
\definecolor{rwthblue}{HTML}{00549F}
\definecolor{rwthblue25}{HTML}{C7DDF2}
\definecolor{rule}{HTML}{8EBAE5}

% --- header banner (typographic, no logo asset) -----------------------------
\newcommand{\rwthbanner}{%
  \noindent\colorbox{rwthblue}{%
    \parbox{\dimexpr\textwidth-2\fboxsep\relax}{%
      \vspace{2.5pt}%
      {\color{white}\textbf{RWTH AACHEN UNIVERSITY}}\hfill
      {\color{white}\footnotesize Computational Engineering Science}%
      \vspace{2.5pt}}}%
  \vspace{1pt}\par\noindent{\color{rule}\rule{\textwidth}{1.3pt}}%
}

\begin{document}

% ---- title block -----------------------------------------------------------
\rwthbanner
\vspace{0.8em}

{\raggedright
  {\LARGE\bfseries Multiphysics Hub}\\[0.35em]
  {\large A real-time browser simulation in which one car couples\\
   vehicle dynamics, aerodynamics, and vehicle--bridge interaction}\\[0.8em]
  {\normalsize Muhammet Emir Fil}\\[0.15em]
  {\small B.Sc.\ Computational Engineering Science, RWTH Aachen \quad\textbullet\quad Technical report, \today}\par
}
\vspace{0.6em}
\noindent{\color{rule}\rule{\textwidth}{0.6pt}}
\vspace{1.0em}

% ---- AIM (before the abstract, as requested) -------------------------------
\noindent\colorbox{rwthblue25}{%
  \parbox{\dimexpr\textwidth-2\fboxsep\relax}{%
  \vspace{3pt}\textbf{Aim.} This project pulls three otherwise separate simulation
  studies, vehicle dynamics, aerodynamics, and vehicle--bridge interaction, into one
  real-time application in which a single car couples all three. The wider goal is to
  show two things at once: that coupled multiphysics can run interactively in a plain
  web browser, and that such an interactive demo can still be tied to a verified
  numerical reference instead of drifting into a look-alike animation. Each model is
  built and checked against ground truth in Python, then re-implemented for the browser
  and compared back against that Python reference.\vspace{3pt}}}

\vspace{1.0em}

% ---- abstract --------------------------------------------------------------
\begin{abstract}
\noindent
A car is driven in real time around a user-drawn track. As it moves it couples three
physics models that are recomputed every frame, with no pre-recorded data: a double-track
(four-wheel) vehicle with Pacejka tyres and load transfer, a 2D fluid solver around the
car silhouette that returns a drag and a downforce, and a moving-load bridge beam that
deflects under the wheel force. Each model is first written in numpy and checked against
an analytical or finite-element reference, then ported to JavaScript for the browser and
compared against its numpy version on identical inputs (agreement of order $10^{-13}$ for
the vehicle and $10^{-9}\,$m for the bridge). This report describes the three models, states
which couplings are physical and which are only illustrative, and gives the verification
that backs the numbers shown live on screen.
\end{abstract}

% ---------------------------------------------------------------------------
\section{Overview}
% ---------------------------------------------------------------------------
The wider portfolio contains three standalone, separately verified projects: a
minimum-lap-time racing-line optimiser, a vehicle--bridge interaction and Bridge
Weigh-in-Motion simulator, and a Navier--Stokes vortex-street solver. The Multiphysics Hub
connects them through the object they share, a car, into one real-time top-down application.
Driving the car runs all three couplings together (Fig.~\ref{fig:hero}): the tyres generate
the forces, the air flowing past the body adds drag and downforce, and crossing a bridge
deflects it. The point of the exercise is to keep an interactive, responsive program honest:
the live readouts are the verified models, not a stand-in.

\begin{figure}[h]
  \centering
  \includegraphics[width=0.95\textwidth]{hub_bridge_shot.png}
  \caption{The three physics running together. The car crosses the bridge (centre). The aero
  panel (top right) shows the 2D wake with the current drag and downforce; the bridge panel
  (bottom) shows the 40\,m span deflecting under the moving load, coloured by bending moment,
  with $f_1 = 2.60\,$Hz and the dynamic amplification factor; the left panels show the
  friction-circle state and the per-wheel load transfer. Every value is computed in the
  current frame.}
  \label{fig:hero}
\end{figure}

% ---------------------------------------------------------------------------
\section{The three models}
% ---------------------------------------------------------------------------

\subsection{Vehicle: a double-track car with Pacejka tyres}
The car is a planar double-track (four-wheel) model. Each tyre's lateral force follows the
Pacejka Magic Formula,
\begin{equation}
  F_y = D\,\sin\!\Big(C\,\arctan\!\big(B\alpha - E(B\alpha - \arctan B\alpha)\big)\Big),
  \qquad D = \mu_\text{eff}(F_z)\,F_z,
\end{equation}
and the longitudinal force is friction-ellipse coupled,
$F_x = \mathrm{thr}\cdot\sqrt{\max\!\big(0,(\mu F_z)^2 - F_y^2\big)}$, so a tyre cannot give
full grip sideways and lengthways at the same time. The body-frame Newton--Euler equations,
including the centripetal $v_y r$ and $v_x r$ terms, integrate the state
$[X,Y,\psi,v_x,v_y,r]$ with RK4. The four wheels carry longitudinal and lateral load transfer,
and since the tyre's peak friction falls as load rises ($\mu_\text{eff}(F_z)$ concave),
transfer changes each axle's total grip. That is what produces the car's understeer or
oversteer balance, set here by the front load-transfer share.

\subsection{Aerodynamics: a 2D fluid around the car}
A 2D incompressible fluid solver runs around the car silhouette every frame, using Stam's
stable-fluids scheme (semi-Lagrangian advection with a Gauss--Seidel pressure projection) and
imposing the body by volume penalisation, the same immersed-boundary idea as in the
portfolio's Navier--Stokes solver. It draws the wake and, through the standard coefficient
relation $F = \tfrac12\rho C A V^2$, returns a drag and a downforce that raise the tyre loads;
at $60\,$m/s the downforce is about $30\,\%$ of the car's weight and lifts the rear-grip cap.
This panel is labelled qualitative on screen: it is a 2D, model-Reynolds field that shows
separation and the wake, not a quantitative aerodynamic prediction. The benchmark-verified
staggered-grid solver is the separate vortex-street project, which the report cites for that
purpose rather than claiming it here.

\subsection{Bridge: a moving-load modal beam}
The bridge is a modal reduction of the Euler--Bernoulli moving-load beam from the VBI project:
the first few mode shapes carry the response, integrated under the travelling wheel force. It
deflects live under the car, coloured by bending moment, and reports the dynamic amplification
factor. Two couplings are kept physical and quantitative: aerodynamic downforce changes the
tyre grip, and the wheel load (including that downforce) deflects the span.

% ---------------------------------------------------------------------------
\section{Verification}
% ---------------------------------------------------------------------------
The interactivity should not cost correctness, so verification runs on two layers. The numpy
models are the reference and are checked against ground truth; the JavaScript ports that run in
the browser are then checked against those numpy models on identical inputs. Both integrate at
a fixed \SI{250}{\hertz} with RK4, decoupled from the render rate.

\paragraph{Layer 1: each model against a reference.} The vehicle is checked against static and
steady-state references; the bridge against the closed-form static deflection and against the
independent full Euler--Bernoulli FEM eigensolve from the VBI project; the fluid solver against
its own conservation properties (Table~\ref{tab:verify}).

\begin{table}[h]
  \centering
  \caption{Each model checked against an analytical or FEM reference.}
  \label{tab:verify}
  \footnotesize
  \setlength{\tabcolsep}{5pt}
  \begin{tabular}{@{}lll@{}}
    \toprule
    Check & Reference & Result \\
    \midrule
    Static wheel loads (at rest)          & $mg/4$ each                          & exact \\
    Vertical equilibrium (any $a_x,a_y$)  & $\sum F_z = mg$                      & exact \\
    Longitudinal / lateral load transfer  & $m|a_x|h/L$,\; $2 s\,m\,a_y h/\text{tr}$ & match closed form \\
    Low-speed cornering                   & Ackermann $\delta \approx L/R$       & match \\
    Balance trend                         & more front transfer $\Rightarrow$ wider $R$ & monotonic \\
    Bridge static mid-span deflection     & $PL^3/48EI$                          & rel.\ err $< 10^{-3}$ \\
    Bridge modal frequencies              & full VBI Euler--Bernoulli FEM        & within \SI{2}{\percent} ($f_1=2.60$\,Hz) \\
    Bridge DAF (slow / at speed)          & $\to 1$ quasi-static; bounded $>1$   & confirmed \\
    Fluid solver                          & divergence drops; stable; sheds a wake & confirmed \\
    \bottomrule
  \end{tabular}
\end{table}

\paragraph{Layer 2: JavaScript against Python.} Each browser model is run against its numpy
twin on the same inputs (Fig.~\ref{fig:checks}). The vehicle port agrees with Python to order
$10^{-13}$ over thousands of steps; the bridge crossing agrees to order $10^{-9}\,$m against a
$0.23\,$mm peak deflection, the residual being transcendental round-off accumulated over the
ring-down. This is the basis for trusting the live numbers: they are the verified model rather
than a visual approximation of it.

\begin{figure}[h]
  \centering
  \begin{subfigure}{0.48\textwidth}
    \includegraphics[width=\linewidth]{xcheck_shot.png}
    \caption{JavaScript vs Python on identical inputs.}
  \end{subfigure}\hfill
  \begin{subfigure}{0.48\textwidth}
    \includegraphics[width=\linewidth]{aero_test_shot.png}
    \caption{Fluid-solver self-check (divergence, stability, wake).}
  \end{subfigure}
  \caption{Layer-2 verification: the shipped JavaScript is held to the numpy reference.}
  \label{fig:checks}
\end{figure}

% ---------------------------------------------------------------------------
\section{The application}
% ---------------------------------------------------------------------------
The page (Fig.~\ref{fig:hub}) opens in a self-driving demo loop and hands control to the user
on any drive key. It is a top-down circuit with a live minimap, a friction-circle readout,
per-wheel load bars that grow and redden under transfer, a manual gearbox, handbrake and
reverse, and an editable track. It runs entirely client-side: open \texttt{index.html}, with no
server and no build step, and it is deployed on GitHub Pages.

\begin{figure}[h]
  \centering
  \includegraphics[width=0.74\textwidth]{hub_shot.png}
  \caption{The driving view: friction-circle state, per-wheel load transfer, inputs and state,
  and the minimap, all updated each frame from the verified vehicle model.}
  \label{fig:hub}
\end{figure}

% ---------------------------------------------------------------------------
\section{Limitations}
% ---------------------------------------------------------------------------
\begin{itemize}
  \item The whole model is two-dimensional.
  \item The aero field is qualitative (2D, model-Reynolds). It drives the car through standard
        coefficient relations, not through a quantitative pressure integration; the
        benchmark-verified solver is the separate vortex-street project.
  \item The bridge is a modal reduction of a simply-supported span, not a continuous multi-span
        girder. The gearbox is a driver-input layer on top of the vehicle dynamics.
  \item Coupling is one-way where it is only illustrative (the aero render) and two-way only
        where it is physical (downforce to grip, wheel load to deflection).
\end{itemize}

% ---------------------------------------------------------------------------
\section{Reproducibility}
% ---------------------------------------------------------------------------
\begin{sloppypar}
The models are numpy-only and re-implemented in JavaScript. The checks ship with the code:
\texttt{verify\_car\_dynamic.py} and \texttt{verify\_doubletrack.py} for the vehicle,
\texttt{verify\_bridge.py} for the beam (against $PL^3/48EI$ and the full VBI FEM),
\texttt{web/xcheck.html} for JavaScript versus Python, and \texttt{web/aero\_test.html} for the
fluid-solver self-check. The application is \texttt{web/index.html}; the live build runs on
GitHub Pages.
\end{sloppypar}

% ---------------------------------------------------------------------------
\begin{thebibliography}{9}
\small
\bibitem{pacejka} H.\,B.\ Pacejka, \emph{Tyre and Vehicle Dynamics}, 3rd ed.,
  Butterworth-Heinemann, 2012. (Magic Formula tyre model.)
\bibitem{milliken} W.\,F.\ Milliken \& D.\,L.\ Milliken, \emph{Race Car Vehicle Dynamics},
  SAE International, 1995. (Load transfer and handling balance.)
\bibitem{stam} J.\ Stam, ``Stable fluids,'' \emph{Proc.\ SIGGRAPH}, 1999.
  (Semi-Lagrangian advection with pressure projection.)
\bibitem{angot} P.\ Angot, C.-H.\ Bruneau \& P.\ Fabrie, ``A penalization method to take into
  account obstacles in incompressible viscous flows,'' \emph{Numerische Mathematik}, 81, 1999.
\bibitem{clough} R.\,W.\ Clough \& J.\ Penzien, \emph{Dynamics of Structures}, 3rd ed.,
  Computers \& Structures, 2003. (Modal superposition.)
\bibitem{fryba} L.\ Frýba, \emph{Vibration of Solids and Structures under Moving Loads},
  Noordhoff, 1972. (Moving-load reference solution.)
\end{thebibliography}

\end{document}
